% SPDX-License-Identifier: CC-BY-4.0
% Copyright (c) 2026 Martin Schr\"oder <info@swedishembedded.com>

\section{Evaluation protocol}
\label{sec:protocol}

The protocol below was fixed before any comparison run was trained. Anything
not listed here is exploratory and is labelled so where it appears.

\paragraph{Corpus and metric.} All models are trained on the TinyStories
training split \citep{eldan2023tinystories} and evaluated on its validation
split, which no model or tokenizer sees during training. The primary metric is
bits per byte (BPB) on the validation split: the summed next-token negative
log-likelihood in bits divided by the UTF-8 byte count of the evaluated text.
BPB does not depend on the tokenizer, so models with different vocabularies,
including the pretrained Qwen3-0.6B, are scored on the same text.

\paragraph{Models.} Every trained model is a pre-norm decoder with the
attention sublayer of \cref{sec:method}; they differ only in width and in the
feed-forward sublayer.
\begin{itemize}
\item \textbf{SpectralLM} $\mathcal{S}(m,D,C)$: residual width $m$, virtual
  signal length $D$, $C$ mixer channels.
\item \textbf{Narrow control} $\mathcal{Q}(m,f)$: the same decoder with a
  SwiGLU feed-forward network of hidden width $f$, at the same residual width
  $m$, same attention shape, same embedding and head. By
  \cref{prop:absorption} this is also the ``truncated spectrum with dense
  maps'' model. It is run at two hidden widths: $f_{\mathrm{std}}$, a
  conventional ratio, and $f_{\mathrm{par}}$, chosen so that its non-embedding
  parameter count matches $\mathcal{S}$.
\item \textbf{Wide control} $\mathcal{Q}(D,f)$: a conventional decoder at the
  virtual signal width $D$, the model a spectral state of length $D$ would be
  compressing.
\item \textbf{Pretrained reference}: Qwen3-0.6B \citep{qwen3report}, scored
  zero-shot on the validation split with its own tokenizer, and after
  fine-tuning on the training split for the token budget of the trained models.
\end{itemize}

\paragraph{Training.} All trained models use the same tokenizer, context
length, batch size, token budget, AdamW settings, cosine schedule and
precision (fp32). The learning rate is chosen per model family from the same
grid on a short run, scored on a slice of the training split, and fixed before
the comparison runs. Each comparison configuration is trained with three
seeds drawn at random and recorded.

\paragraph{Hypotheses.}
\begin{description}
\item[H1 (primary).] At equal residual width, equal attention, and equal
  training tokens, $\mathcal{S}(m,D,C)$ reaches a mean validation BPB no
  worse than $\mathcal{Q}(m,f_{\mathrm{par}})$ by more than $0.02$ BPB. It is
  rejected if the mean difference over seeds exceeds $0.02$ BPB.
\item[H2.] $\mathcal{S}(m,D,C)$ does not reach the quality of
  $\mathcal{Q}(m,f_{\mathrm{std}})$, which has more parameters; the gap is
  reported with its uncertainty either way.
\item[H3.] On Qwen3-0.6B's residual stream, the lowest $m$ modes of a real
  Fourier basis retain no more energy, and no more next-token information, than
  $m$ random orthogonal directions, and less than the top $m$ principal
  components.
\item[H4.] The fast-transform mixer and the dense-product mixer compute the
  same function within single-precision tolerance, and the fast one is faster
  at every evaluated $(m,D,C)$.
\end{description}

\paragraph{Correctness gates.} Before any comparison run: the basis
satisfies orthonormality and Parseval's identity to rounding error; every new
operator and the whole model pass a finite-difference gradient check; the CPU
and GPU backends agree; and the model drives the loss of a fixed small batch
to near zero.

\paragraph{Reporting.} For every configuration we report the mean and range
over seeds of validation BPB, the non-embedding and embedding parameter
counts, the analytic arithmetic per token, measured training throughput, and
peak device memory. Every number is generated from a result record that names
the commit, command, seeds and device.
